\section{The model}\label{sec:model}
All calculations hold the baryonic mass fixed at $0.1\Ms$ and omit rotation,
magnetism, mass loss, accretion and external heating. The initial hydrogen
mass fraction is $X=0.7$ and the metal fraction $Z=0.02$, with the relative
metal abundances of the Grevesse--Sauval (GS98) mixture and the remainder in
$^4$He. Tracks S, F and P start with no $^3$He; the continuous calculation
starts with initial deuterium. Table~\ref{tab:ingredients} distinguishes
the physical assumptions of the comparison tracks from those selected for
the continuous calculation.

{\small\renewcommand{\arraystretch}{1.12}
\begin{longtable}{@{}>{\raggedright\arraybackslash}p{0.13\textwidth}>{\raggedright\arraybackslash}p{0.44\textwidth}>{\raggedright\arraybackslash}p{0.37\textwidth}@{}}
\caption{Physical ingredients by calculation.}\label{tab:ingredients}\\
\toprule Component & Tracks S, F, P and flash continuations & Continuous calculation \\\midrule\endfirsthead
\toprule Component & Tracks S, F, P and flash continuations & Continuous calculation \\\midrule\endhead
Starting model & Static, fully convective hydrogen-burning model at age zero:
$R=0.1259\Rs$, $\Teff=2766$ K, $L=0.0008362\Ls$, central density
381.7 g cm$^{-3}$. & Fully convective Hayashi-track model: $\Teff=2864$ K,
$R=1.594\Rs$, central temperature 0.5398 MK, with initial deuterium. \\[4pt]
Nuclear reactions & Reduced ppI/ppII network with $^3$He out of equilibrium,
Solar Fusion II rates \cite{adelberger11} and finite-degeneracy Salpeter--Van
Horn screening. Tracks F and S add a time-dependent CN cycle ($^{12}$C,
$^{13}$C, $^{14}$N) with the $^{14}$N capture rate of Solar Fusion III
\cite{acharya25}; in Track S nuclear captures change the metal mass. Track P
has no CN burning. Pep, hep, ppIII and the oxygen branches are absent in all.
& Initial deuterium, pp and closed CN burning with Solar Fusion III rates
for both networks and the same screening prescription.
Pep, hep, ppIII and oxygen branches absent. \\[4pt]
Heat transport & Mixing-length convection with the Ledoux criterion;
radiative opacity plus electron conduction from Ioffe/Cassisi mixtures;
plasma-neutrino losses \cite{hrw94}. Other thermal-neutrino processes are
omitted. & Ledoux mixing-length convection, radiation, plasma-neutrino
losses and conductive heat transport. In hot layers the screened kinetic
calculation supplies the conductive coefficient; it joins the cool-envelope
Ioffe/Cassisi coefficient smoothly between 2 and 3 MK. \\[4pt]
Composition transport & Track P: instantaneous mixing within each convective
region, slow mixing in marginally stable layers, no microscopic diffusion.
Track F: microscopic diffusion of H, $^3$He, C and N with the metal
inventory fixed; chemical forces from the equation-of-state free energy;
electron collisions averaged over the finite-temperature Fermi distribution.
Track S: as F, with all heavy elements moving as one group with distinct
charges, and the equation of state and opacity depending on the total metal
abundance. Flash continuations: finite mixing $D=v\ell/3$ in hot convective
regions; semiconvective and thermohaline coefficients set to zero.
& Instantaneous whole-star mixing while the star is fully convective, with
the heat carried by changing composition included after deuterium
exhaustion. Screened H/He/metal species fluxes act across radiative
boundaries at temperatures of at least 2 MK; the metal law assumes full
ionization. Cooler radiative species transport is unsupported.
The heat law, rather than the species flux, uses the smooth 2--3 MK join. \\[4pt]
Equation of state & FreeEOS 3.0 free-energy tables with numerical electron
integrals, interpolated through a shared Helmholtz potential so that
pressure, energy and composition responses stay consistent. Track P: 72
composition tables. Track F: 152 tables at 38 hydrogen and four $^3$He
fractions. Track S: a variable-metal family from $Z=0$ to 0.08, extended to
0.16 for the metal-enriched core. None describes the cold, dense remnant.
& One interpolated potential with planes at $Z=0$, 0.005, 0.02, 0.04 and
0.16, twelve hydrogen-share coordinates and four $^3$He-share coordinates;
all 512 zones support the required composition derivatives. \\[4pt]
Opacity & AESOPUS 2.1 gas opacity at low temperature (no grains) and Los
Alamos ATOMIC/TOPS opacity at high temperature \cite{colgan16,tops},
including hydrogen-free mixtures above $5.012\times10^5$ K. Composition
continuations in $X$ and $Z$ are given in Appendix~\ref{app:tech}.
& Same tables and continuations. \\[4pt]
Atmosphere & TLUSTY/SYNSPEC plane-parallel LTE atmospheres with
frequency-resolved absorption, scattering and collision-induced absorption,
no condensates, joined to the interior at Rosseland optical depth 100.
Track P: 220 atmospheres over $X=0.1$--0.7, $^3$He 0--0.12, through
5600 K. Track S: additional families down to $Z=10^{-8}$ and $1-X=10^{-5}$,
4600--6000 K, $\log g$ to 5.95, with a local continuation to 6.1. Flash
continuations: the pure-hydrogen white-dwarf table distributed with MESA
\cite{rohrmann11,mesa_atm}, joined at optical depth 25.12 and extrapolated
below $\log g=5.5$. & A composition-dependent family at $X\geq0.65$,
$^3$He $\leq0.0075$, fixed $Z=0.02$, joined at optical depth 100. The
continuous metal-dependent selection joins measured families through
nearly pure hydrogen, including trace helium, at the same optical depth.
The 4600 K gravity row and its high-gravity corner have solved source
columns; intermediate knots interpolate the sources. Trace helium remains
approximated. Coverage ends at 4600 K and $\log g=6.1$. \\[4pt]
Mesh and time integration & 512 mass points, fixed for Tracks S, F and P;
flash grids refined locally to 578--2115 points. Structure, burning and
mixing are solved together implicitly. Every accepted step is compared with
two half steps; species inventories, inert-metal mass, nuclear rest-mass
release and the first law are checked independently. Saved models restart
exactly. & 512 zones; the same step-doubling check; independent isotope,
mass-defect and first-law audits on every interval; exact restart replay.
Radiative-core evolution and restart are registered tests. \\
\bottomrule
\end{longtable}}

For Tracks S, F and P the relative time-error tolerances are 0.0003 for radius, density and
temperature, 0.003 for local $^3$He and 0.001 for the other species; the
abundance solver converges to $10^{-10}$. Tightening the structure and
$^3$He tolerances threefold over a matched 5 Gyr interval changes luminosity
by 0.002872\% and surface temperature by 0.01705 K, while loosening the
structure tolerance to 0.001 over 2 Gyr changes luminosity by 1.078\% and
temperature by 3.114 K; the tracks therefore keep the original settings.
The continuous calculation uses a structure tolerance of $10^{-4}$ and an
energy tolerance of 0.005. Local composition solves converge to $10^{-15}$
in a single instantaneously mixed region and $10^{-12}$ in a stratified
star; the separate global abundance-balance tolerance remains $10^{-14}$.

Throughout the paper a \emph{paired comparison} means two calculations
started from the same model, differing in one setting, and compared at the
same age. Such comparisons measure sensitivity to that setting over the
stated interval; they are not error bounds for a complete lifetime, and the
text says so wherever the distinction matters.

Replacing the Solar Fusion II rates by Solar Fusion III with everything else
fixed raises the initial luminosity by 0.7454\% and the radius by 0.2164\%.
FreeEOS omits potassium (mass fraction $4.56\times10^{-6}$). The interior
equation of state, atmosphere chemistry, opacity and conduction come from
different sources; their mutual consistency is checked at sampled states
(Appendix~\ref{app:tech}) rather than guaranteed.
